% 06-verifiers.tex: независимая проверка и первые свидетели извне.
\section{Проверка без эмитента: пакетные верификаторы, контракт соответствия и первые посторонние}
\label{sec:verifiers}

\subsection{Отделяемый верификатор как вещь, которую устанавливает посторонний}

\code{polaris-verify} проверяет пакет подлинности удостоверения одной лишь стандартной библиотекой
ML-DSA: ни строки кода Polaris, ни базы данных, ни сети, если только его не попросили запросить
состояние. Версия 2 описывала его как сценарий в репозитории. Теперь это пакет в публичном
реестре, на \code{1.0.0-rc.3}, кандидат в выпуски, под которым с 15 сентября 2026 года лежит
\code{0.1.0}, опубликованный собственным рабочим процессом репозитория по доверенной публикации,
так что нигде не существует долгоживущих учётных данных, а сценарий, который называла версия 2,
есть переходник на него. Он отказывается запускаться, пока запуск не скажет, какую криптографию
выполняет: \code{--pqc-provider oqs}, \code{cryptography} или \code{auto}, либо
\code{--dev-placeholder}, без значения по умолчанию и без переменной окружения, потому что
верификатор, тихо переставший проверять, всё равно отвечает на каждый вопрос. Всякий
машиночитаемый вердикт несёт поле \code{crypto}, а под флагом разработки настоящие подсистемы
выключаются, так что запуск для разработки вообще не способен произвести истинный вердикт;
построение этого запора обнаружило, что первая версия флага объявляла режим, в него не входя, а это
ярлык, а не режим. \sImpl

Это единственная составная часть Polaris, вход которой выбирает противник, потому что доверяющая
сторона запускает её на том, что предъявил посторонний, и это та часть, через которую проверяется
всякий последующий предмет протокола: описи, контрольные точки эпох, ленты отзывов, наборы
состояний, утверждения о состоянии, вершины прозрачности, расписки, отметки времени, подписанные
документы, реестры, списки доверия, предъявления, доказательства держателя, полномочия агентов.
Три независимые реализации ML-DSA сходятся на опубликованных векторах: решёточная библиотека,
OpenSSL и третья в TypeScript-SDK. \emph{Тотальность} против враждебного входа удерживается
метаморфическим фаззером, а рядом с ним работает набор атак из исполняемых противников; grep
доказывает, что строка существует, атака доказывает, что защита держит. \sImpl

Способен ли посторонний в действительности его установить, измеряется, а не предполагается. При
каждой отправке и перед каждой публикацией учение собирает колесо, ставит его в одноразовое
окружение, где нет ничего от Polaris, и запускает из каталога, который не есть репозиторий, со
снятыми переменными \code{POLARIS\_*}: запретная исполняемая поверхность (веб-приложение, драйвер
базы данных, слой проверок) обязана отсутствовать, команда обязана отказаться запускаться без
объявленного режима криптографии, подлинное удостоверение ML-DSA-87 обязано провериться, а его
подменённый двойник обязан не провериться. Учение несёт отрицательный контроль, второе колесо с
прикрученным к верификатору импортом базы данных, который стенд обязан поймать. Это важно, потому
что всякий иной набор тестов в репозитории работает изнутри репозитория с репозиторием на пути
поиска; импорт, который остался бы там зелёным, падает на машине первого же постороннего. \sImpl

\subsection{SDK и контракт соответствия}

Доверяющая сторона ставит один из двух проверочных SDK, \code{polaris-sdk-python} и
\code{polaris-sdk-ts}, каждый из которых отвечает на один узкий вопрос о предъявленном
удостоверении, подлинно ли оно и полномочно ли прямо сейчас, и никогда не возвращает данных
человека. Оба лежат в своих реестрах. Тот, что на TypeScript, есть урок о том, чему учит реестр:
его первая упаковка направляла \code{exports} на исходный код TypeScript, который Node внутри
\code{node\_modules} обрабатывать отказывается, так что пакет не смог бы импортировать ни один
потребитель ни на одной версии Node; тесты репозитория работали внутри каталога пакета, где
обработка разрешена, и ничего не заметили. Теперь он компилируется в дистрибутив, а учение ставит
упакованный архив в пустой проект. \sImpl

Оба удерживаются языконезависимым набором проверок соответствия: исполнитель гоняет любой
верификатор через стандартный ввод и вывод по опубликованным случаям, коих на \code{1.0.0-rc.7}
насчитывается 118 по двадцати пяти типам предметов, включая составное решение о доверии в
федерации, доказательство держателя, полномочие агента и его отзыв. Прохождение этого набора и
означает \emph{соответствие}, а нормативная спецификация формата передачи (язык RFC 2119, двадцать
семь форматных строк, точный перечень подписываемых полей каждого предмета) есть то, против чего
строила бы реализация, не импортирующая ни строки кода Polaris. \sImpl

Затем у контракта спросили, чего он не спрашивает, измерением, и ответ ограничивает то, что
означает соответствие. У трёх типов предметов не было ни одного случая с испорченной подписью, так
что верификатор, ни разу их подпись не проверявший, соответствовал. Пять из опубликованных
действительных векторов истекли месяцами раньше, и ни один случай этого не заметил, потому что ни
один не утверждал свежести. Вопрос о свежести вскрыл два расхождения между поставляемыми
эталонными верификаторами: ни один не ограничивал возраст доказательства держателя, так что
интегратор, следовавший за одним из них, не получал защиты от повторов при предъявлении, и ни один
не умел сообщить, доверенный ли у предмета эмитент вообще, так что подлинная подпись постороннего
проходила. Всё исправлено, и три верификатора сходятся на каждом случае. Чего контракт по-прежнему
не доказывает, изложено там, где интегратор это прочтёт: 43 поля вердикта не стеснены ни одним
опубликованным случаем, и ни одно из них нельзя закрыть, написав такой случай, потому что каждому
нужен соответствующий верификатор, вычисляющий то, чего он не вычисляет. Контракт стесняет то, что
вычисляет его слабейшая соответствующая реализация. Затем мутировали и сами эталонные
верификаторы: каждый отказ в каждом из них поочерёдно обращали, и восемнадцать из восемнадцати в
реализации на Python и девять из четырнадцати в реализации на TypeScript приняли то, ради отклонения
чего они и существуют, при зелёных тестах самого SDK и зелёном исполнителе соответствия, включая
исчерпанный счётчик использований в полномочии агента и, в TypeScript-SDK, сравнение за постоянное
время, у которого обращённая проверка длины заставила пятибайтовое значение сравняться с
32-байтовым корнем. Всё это закрыто в собственных тестах SDK, а учение обращает все 32 отказа при
каждой отправке с отдельным отрицательным контролем на каждый SDK. \sImpl

\subsection{Версионирование и замороженная версия 1}

Старшие версии протокола живут в форматной строке (\code{polaris-registry/1}), младшие
объявляются реестром и для проверки никогда не нужны, а правило согласования нормативно: отвергни
неизвестную старшую, прими любую младшую, никогда не выпускай необъявленной версии, ответь на
неподдерживаемую определённой ошибкой и списком поддерживаемых. Версия 1 протокола заморожена под
контрольными суммами, рядом с ней помещён отделяемый верификатор \code{v9.317}, а набор тестов на
совместимость работает в обе стороны при каждой отправке: сегодняшние верификаторы обязаны
принимать всё, что опубликовала версия 1 (44 случая из 44 держатся), а закреплённый старый
верификатор обязан никогда не принимать того, что сегодняшний набор отвергает. \sImpl

% Рисунок открывает подраздел под его заголовком, и заголовок с рисунком остаются на одной
% странице: нужное им вместе место резервируется до заголовка, так что страница, которая не
% вмещает обоих, разрывается перед заголовком, а не между ними.
\Needspace{36\baselineskip}
\subsection{Снаружи: верификатор для кошельков, которых Polaris никогда не встречал}

\figfile{outside-witnesses}{Три вещи за пределами репозитория, которые его задействовали, и два
контроля, без которых первая из них ничего бы не значила. Кошелёк, написанный здесь никем,
предъявил удостоверение и был принят; тот же кошелёк, когда верификатор доверял другому ключу
эмитента, получил отказ; и повторное предъявление против уже отвеченного запроса получило отказ.
Размещённый пакет Фонда сам засчитал семь отказов и оставил четыре принятия на человека-рецензента.
Опубликованный набор векторов сходится с обоими свидетелями подписи на примитиве. Ни одна из трёх
не есть аудит, сертификация или пилотное внедрение.}{fig:outside-witnesses}

Всё вышеизложенное есть Polaris, проверяющий Polaris, и честной границей версии 2 было то, что ни
одна внешняя реализация с ним не взаимодействовала. Теперь это предложение иное, на одном пути, а
путь был выбран потому, что на нём в действительности и говорит сторонний кошелёк. По операционному
договору первой целью по совместимости был верификатор OpenID4VP 1.0 по профилю высокой
надёжности взаимодействия, а у пакета проверок соответствия OpenID Foundation есть план тестов
ровно для этой роли, \code{oid4vp-1final-verifier-haip-test-plan}, в котором \emph{удостоверение
подписывает сам пакет}, так что собственная выдача Polaris под ML-DSA в этом не участвует.
\code{polaris-oid4vp} есть этот верификатор: он проверяет SD-JWT VC, подписанный ES256, с
привязкой к ключу держателя, доставленный как зашифрованный ответ \code{direct\_post.jwt}, и
отказывает с названной причиной, по одной на каждый отрицательный модуль плана:
\code{issuer\_signature}, \code{sd\_hash}, \code{kb\_signature}, \code{nonce}, \code{audience},
\code{kb\_freshness}. Это отдельный пакет, потому что \code{polaris-verify} обещает не открывать
ни одного сокета, а этому слушать необходимо. Одна из его заготовок не замкнута на себя: полный
ответ, произведённый собственным кошельком пакета проверок и построенный Java-библиотекой, которой
этот проект не пользуется, помещён в дерево, расшифровывается и проверяется, а годом позже
получает отказ под другим одноразовым числом, под другой аудиторией и с одним переписанным
раскрытием. \sImpl

Затем произошли две вещи, которых репозиторий не делал сам с собой. 15 сентября 2026 года
размещённая служба проверки соответствия OpenID Foundation прогнала все одиннадцать модулей плана
против верификатора через открытый интернет, забирая объект запроса и отправляя ответ снаружи:
ноль отказов, ноль предупреждений. Семь отрицательных модулей, в которых кошелёк пакета посылает
предъявление, сломанное одним определённым образом, несут собственный \code{result: PASSED}
службы, засчитанный автоматически по отказу верификатора, что делает их самым близким к внешнему
противнику, что этому проекту довелось получить. Четыре положительных модуля завершились в
состоянии \code{REVIEW}, которое пакет отводит человеку, чтобы тот удостоверил, что верификатор
отобразил успешную проверку; снимки экрана приложены, а рецензент Фонда видит их лишь при подаче
на сертификацию, которая не запрашивалась. REVIEW не есть PASSED, и Polaris не сертифицирован.
Первое описание прогона говорило: одиннадцать из одиннадцати чисто, и это было ложью в лестную
сторону: для каждого модуля использовался один закреплённый псевдоним, а пакет прерывает тест,
когда под тем же псевдонимом начинается другой, и четыре положительных были убиты своими
преемниками, что пакет уже записал внутри того самого журнала, который скачали ради подтверждения
утверждения. Прогон повторили по одному модулю за раз. Прогон несёт два отрицательных контроля, по
одному на направление: верификатор, пропатченный отказывать во всём, замечается положительными
модулями, а пропатченный принимать всё, всеми семью отрицательными. \sWit{} для профиля; \sExt{}
для сертификации.

И в тот же день неизменённый walt.id Wallet API v2, опубликованный образ по закреплённому
дайджесту, предъявил SD-JWT VC верификатору по OpenID4VP и был принят: он забрал подписанный
объект запроса, сверил его с зарегистрированным якорем доверия, сопоставил удостоверение с
запросом, подписал токен привязки к ключу ключом P-256, который выработал сам и которым этот
репозиторий никогда не владел, зашифровал ответ, а \code{polaris-oid4vp} проверил всю цепочку.
Сперва он Polaris отверг, и был прав. Собственный генератор ключей верификатора произвёл
сертификат для подписи запросов вовсе без расширения назначения ключа, так что сертификат, вся
работа которого состоит в подписании объектов запроса, не был помечен как пригодный для подписи;
одиннадцать модулей соответствия, 143 теста пакета, мутационное учение по каждому отказу и каждая
проверка в слое инвариантов прошли мимо этого, а первая же реализация извне отказалась
продолжать. Исправлено в тот же день, с тестом, который краснеет на листе, где расширения нет, и
засчитано в сводной таблице как единственное исправление, вызванное внешней находкой. Обмен
держится на тех же двух контролях, потому что верификатор, принимающий всё, печатает ту же самую
строку об успехе. Чего он не устанавливает, изложено там же: один кошелёк, один формат
удостоверения, один путь предъявления, ES256 повсюду, что и закрепляет профиль; ничего о любом
ином кошельке, об ISO 18013-5 или о постквантовом пути, и никаких утверждений о том, что
верификатор совместим вообще. Всякий может повторить этот обмен с чистой машины примерно за десять
минут, не клонируя репозитория; страницу, объясняющую как, исполнили прежде, чем написали.
\sWit{} для этого пути.

Родной протокол, форматы \code{polaris-*/1}, на которых построено всё остальное в настоящем
документе, остаётся там, где его оставила версия 2: спецификация, случаи и два SDK существуют, а
реализации сторонней командой по одним лишь документам не случилось. \sExt

Два формата-моста делают границу явной. Удостоверение Polaris отображается в структуру
мобильного документа ISO/IEC 18013-5, и считыватель, говорящий на этом стандарте, его разбирает,
обходит его пространства имён и сверяет дайджест каждого раскрытого элемента с подписанным
объектом безопасности, причём доказывает это независимая реализация CBOR, а не та, что писала
байты; подпись эмитента считыватель проверить не может, потому что это ML-DSA-65, а стандарт
предписывает классические алгоритмы, и репозиторий называет это устройством, а не ограничением,
поскольку постквантовое удостоверение, несущее классическую подпись ради зелёной галочки, было бы
классическим удостоверением. Результат проверки также отображается в проверяемое удостоверение
W3C, намеренно не идентификационное: словарь его субъекта замкнут и отвергает поля личности по
имени, а идентификатор субъекта есть повериферный дескриптор либо отсутствует. Оба суть форматы, а
не модели доверия. \sImpl

\begin{layercard}{Карточка слоя: независимая проверка}
\qa{Что это}{Четыре пакета в публичных реестрах: самостоятельный верификатор, два SDK и верификатор OpenID4VP; исполнитель соответствия по 118 опубликованным случаям; нормативная спецификация формата передачи; замороженный набор версии 1 с межверсионным набором тестов; фаззер, набор атак и три мутационных учения; два формата-моста.}
\qa{Зачем оно существует}{Проверка, выполняемая эмитентом, просит доверяющую сторону довериться серверу эмитента. Подлинность обязана быть разрешима всяким, где угодно, при выключенной сети и программами, которые автор не писал.}
\qa{Чему доверяет}{Опубликованным ключам эмитентов, которым доверяющая сторона решила доверять (её набору якорей); стандартной библиотеке ML-DSA, а на пути OpenID4VP, зарегистрированному якорю доверия X.509 и ES256.}
\qa{Чему отказывает}{Любому коду или базе данных Polaris; любому входу, который не способен разобрать; любому объекту, поданному верификатору не того типа; любой версии, о которой ему не сказали; запуску без объявленного режима криптографии; одноразовому числу или аудитории, которых он сам не запрашивал.}
\qa{Что видит}{Предмет и набор якорей. Для SD-JWT VC, ровно раскрытые утверждения. Больше ничего.}
\qa{Чего не видит}{Сохраняются ли у удостоверения полномочия, если вместе с ним не предъявлен предмет о состоянии; 43 поля вердикта, которых опубликованные случаи не стесняют.}
\qa{Если откажет}{Ошибочно принятая мутация роняет фаззер; отказ, переставший отказывать, роняет мутационное учение SDK; расхождение спецификации и подписанта роняет \code{check\_wire\_spec\_matches\_code}; пакет, который посторонний не может установить, роняет учение о границе; дефект, мимо которого прошёл каждый внутренний инструмент, есть то, ради чего внешнее и нужно.}
\qa{Кем проверяется}{\code{check\_detached\_verifier}, \code{check\_conformance\_suite}, \code{check\_typescript\_sdk}, \code{check\_verifier\_fuzz}, \code{check\_attacks\_run}, \code{check\_protocol\_versioning}, \code{check\_sdk\_refusals\_are\_mutation\_tested}, \code{check\_conformance\_asks\_the\_relying\_party\_question}, \code{check\_oid4vp\_verifier\_boundary}, \code{check\_documented\_verifier\_commands\_run}.}
\qa{Сейчас или потом}{Сейчас, для программ и для одного стороннего кошелька на одном пути. Сертификация, второй кошелёк и реализация родного протокола кем-то ещё суть потом и не в нашей власти.}
\qa{Внутри и снаружи}{Внутри: подписи. Снаружи: слой состояния, отвечающий на вопрос, на который подпись ответить не может.}
\end{layercard}

\nextlayer{Подписание удостоверения доказывает, кто его подписал. Оно не доказывает, сохраняются ли
у этого удостоверения полномочия: печать на отозванном паспорте столь же подлинна, как печать на
действительном. Поэтому Polaris нужно состояние.}
